% image-net post, figure 1: AlexNet (Krizhevsky et al. 2012, Fig. 2), redrawn
% as a flat left-to-right diagram. The two rows are the two GPU streams; every
% box holds that GPU's half of the layer (channel count per GPU). Streams talk
% to each other only at C2->C3 and in the fully-connected layers (crossed
% arrows); everywhere else each GPU sees only its own half.
% Narrow strips after C1, C2 and C5 are response normalisation (fe) and
% overlapping 3x3/stride-2 max pooling (fb), in that order.
% Colour key shared across the post: conv fa, pooling fb, fully-connected fc,
% normalisation fe, softmax/output fd.
\documentclass[tikz,border=3pt]{standalone}
\input{FIGKIT/preamble.tex}
\begin{document}
\begin{tikzpicture}[fig, x=1mm, y=1mm,
    b/.style={tint=#1, rounded corners=1.5pt, line width=.6pt},
    w/.style={->, line width=.5pt, draw=soft, >={Stealth[length=1.5mm,width=1.2mm]}},
    head/.style={font=\rmfamily\small, inner sep=1pt, text height=1.6ex, text depth=.25ex},
    kern/.style={font=\ttfamily\scriptsize, text=soft, inner sep=1pt, align=center},
    size/.style={font=\rmfamily\scriptsize, text=soft, inner sep=1pt},
  ]
  \def\yr{7}     % row centre, GPU 1 at +yr, GPU 2 at -yr
  \def\hh{4.5}   % half box height
  \def\sw{2}     % strip width
  % \pair{name}{left x}{width}{accent}{text}: one box in each GPU row
  \newcommand{\pair}[5]{%
    \foreach \r/\s in {1/1, 2/-1} {
      \draw[b=#4] (#2,\s*\yr-\hh) rectangle ({#2+#3},\s*\yr+\hh);
      \node[font=\rmfamily\footnotesize] at ({#2+#3/2},\s*\yr) {#5};
      \coordinate (#1\r w) at (#2,\s*\yr);
      \coordinate (#1\r e) at ({#2+#3},\s*\yr);
    }
    \coordinate (#1top) at ({#2+#3/2},\yr+\hh);
    \coordinate (#1bot) at ({#2+#3/2},-\yr-\hh);
  }
  % \strips{name}{left x}{accents...}: narrow norm / pool strips after a conv box
  \newcommand{\strip}[4]{%
    \foreach \r/\s in {1/1, 2/-1} {
      \draw[b=#3, line width=.5pt, rounded corners=1pt] (#2,\s*\yr-\hh) rectangle ({#2+\sw},\s*\yr+\hh);
      \coordinate (#1\r e) at ({#2+\sw},\s*\yr);
    }
  }
  % header (name + kernel) above, output size below
  \newcommand{\colhead}[3]{%
    \node[head, anchor=south] (#1h) at ($(#1top)+(0,7.6)$) {#2};
    \node[kern, anchor=north] at (#1h.south) {#3};
  }
  \newcommand{\colfoot}[2]{\node[size, anchor=north] at ($(#1bot)+(0,-1.2)$) {#2};}

  % ---------------- input
  \draw[frame=soft, line width=.6pt, rounded corners=1.5pt] (0,-\yr-\hh) rectangle (7,\yr+\hh);
  \node[font=\rmfamily\footnotesize] at (3.5,0) {RGB};
  \node[head, anchor=south] (inh) at (3.5,\yr+\hh+7.6) {input};
  \node[size, anchor=north] at (3.5,-\yr-\hh-1.2) {$224{\times}224{\times}3$};
  \coordinate (ine) at (7,0);
  % ---------------- C1: 96 kernels 11x11x3, stride 4 -> 55x55x96, norm, pool -> 27x27
  \pair{c1}{12}{8}{fa}{48}
  \strip{c1n}{20.6}{fe}{} \strip{c1p}{23.2}{fb}{}
  \colhead{c1}{C1}{11×11\\stride 4}
  \node[size, anchor=north] at (18.6,-\yr-\hh-1.2) {$55{\times}55{\times}96$};
  \node[size, anchor=north, text=faint] at (18.6,-\yr-\hh-4.6) {pool $\to 27{\times}27$};
  % ---------------- C2: 256 kernels 5x5x48 -> 27x27x256, norm, pool -> 13x13
  \pair{c2}{30.2}{8}{fa}{128}
  \strip{c2n}{38.8}{fe}{} \strip{c2p}{41.4}{fb}{}
  \colhead{c2}{C2}{5×5}
  \node[size, anchor=north] at (36.8,-\yr-\hh-1.2) {$27{\times}27{\times}256$};
  \node[size, anchor=north, text=faint] at (36.8,-\yr-\hh-4.6) {pool $\to 13{\times}13$};
  % ---------------- C3: 384 kernels 3x3x256 (reads both GPUs)
  \pair{c3}{51.4}{8}{fa}{192}
  \colhead{c3}{C3}{3×3}
  \colfoot{c3}{$13{\times}13{\times}384$}
  % ---------------- C4: 384 kernels 3x3x192
  \pair{c4}{65.4}{8}{fa}{192}
  \colhead{c4}{C4}{3×3}
  \colfoot{c4}{$13{\times}13{\times}384$}
  % ---------------- C5: 256 kernels 3x3x192, pool -> 6x6
  \pair{c5}{79.4}{8}{fa}{128}
  \strip{c5p}{88}{fb}{}
  \colhead{c5}{C5}{3×3}
  \node[size, anchor=north] at (84.5,-\yr-\hh-1.2) {$13{\times}13{\times}256$};
  \node[size, anchor=north, text=faint] at (84.5,-\yr-\hh-4.6) {pool $\to 6{\times}6$};
  % ---------------- F1, F2: 4096 each, dropout 0.5
  \pair{f1}{98}{8}{fc}{2048}
  \colhead{f1}{F1}{dropout}
  \colfoot{f1}{$4096$}
  \pair{f2}{114}{8}{fc}{2048}
  \colhead{f2}{F2}{dropout}
  \colfoot{f2}{$4096$}
  % ---------------- output: 1000-way softmax
  \draw[b=fd] (130,-\hh) rectangle (138,\hh);
  \node[font=\rmfamily\footnotesize] at (134,0) {1000};
  \node[head, anchor=south] (outh) at (134,\yr+\hh+7.6) {output};
  \node[kern, anchor=north] at (outh.south) {softmax};
  \coordinate (outw) at (130,0);

  % ---------------- wiring
  \foreach \r in {1,2} {
    \draw[w] (ine) -- (c1\r w);
    \draw[w] (c1p\r e) -- (c2\r w);
    \draw[w] (c3\r e) -- (c4\r w);
    \draw[w] (c4\r e) -- (c5\r w);
    \draw[w] (f2\r e) -- (outw);
  }
  % cross-GPU: every unit reads both halves of the layer below
  \foreach \a/\b in {c2p/c3, c5p/f1, f1/f2} {
    \draw[w] (\a1e) -- (\b1w); \draw[w] (\a2e) -- (\b2w);
    \draw[w] (\a1e) -- (\b2w); \draw[w] (\a2e) -- (\b1w);
  }
  % ---------------- GPU labels
  \node[lab, anchor=east] at (-1.5,\yr) {GPU 1};
  \node[lab, anchor=east] at (-1.5,-\yr) {GPU 2};
  % ---------------- legend
  \begin{scope}[shift={(12,-26)}]
    \foreach \c/\t/\dx/\dy in {fa/{conv + ReLU (channels per GPU)}/0/0, fc/{fully connected + ReLU}/0/-5,
        fb/{max pool 3×3, stride 2 (overlapping)}/52/0, fe/{response normalisation}/52/-5, fd/{softmax}/108/0} {
      \draw[b=\c, line width=.5pt, rounded corners=.8pt] (\dx,\dy-1.4) rectangle ++(2.8,2.8);
      \node[anchor=west, font=\rmfamily\footnotesize, text height=1.6ex, text depth=.4ex] at (\dx+3.8,\dy) {\t};
    }
  \end{scope}
\end{tikzpicture}
\end{document}
